% !TEX root = ../main.tex
\section{Conclusion}

Within the task and evaluation setting studied here, the results support the
use of LLMs as the high-level closed-loop decision core of a
resource-constrained multi-UAV wildfire rescue system. This capability is not
uniformly robust: the complete decision loop remains effective in terms of
basic task completion, but execution quality and mission continuity deteriorate
as operational pressure increases. The proposed LLM-centered closed-loop
decision framework addresses this problem by connecting fire-situation
understanding and resource-constrained mission planning with Validation,
Memory, and Runtime Replanning.

Three sets of evidence support this conclusion. First, the Optimization-Based
Planner remains competitive in the simple and normal episodes, whereas the
complete framework maintains higher basic task completion and execution quality
in the hard and expert episodes across four downstream foundation-model
configurations. This also shows that end-to-end operation is not confined to
one downstream model, although all four configurations lose execution quality
under greater operational pressure. Second, the ablations show condition-dependent
roles: Validation provides an occasional reliability safeguard, Memory becomes
increasingly valuable as resource and scheduling interactions intensify, and
Runtime Replanning is critical when perturbations invalidate an active plan.
Third, Qwen3-8B fine-tuned with validated planning demonstrations retains
98.24\% of the Gemini~3.1~Pro teacher planner's average score on simple and
normal episodes, while reducing mean planning time by approximately 63\% and
planning-token usage by 73\%. This final result supports lower-cost local
deployment of Mission Planning within the evaluated setting.

This answer remains limited to a calibrated simulation, a fixed
episode set, and one real-platform configuration. Moreover, the compact
planner has not yet been evaluated under hard and expert runtime conditions.
The evaluation also does not include a direct numerical comparison with an
external end-to-end rescue method because existing systems assume different
inputs, action spaces, resource constraints, and evaluation protocols. The
reported evidence therefore supports the framework within the task and
evaluation setting defined here. Future work will extend the evaluation to
real-flight and mixed real--sim trials, more diverse platforms and fire
dynamics, and larger fleets, while developing more selective memory retrieval
and more robust, efficient replanning for missions under high operational
pressure.

\endinput
